\subsection*{Fields}
\label{subsec:fields}

\begin{tcolorbox}[title={Definition --- Field}]
\begin{definition}[Field]
\label{def:field}
A \textbf{field} is a set $F$ equipped with two binary operations, addition
$+$ and multiplication $\cdot$, such that:
\begin{enumerate}[label=\textbf{F\arabic*.}]
  \item $(F,+)$ is an abelian group with identity $0$;
  \item $F\setminus\{0\}$ is an abelian group under multiplication with
        identity $1$;
  \item multiplication distributes over addition.
\end{enumerate}
\end{definition}
\end{tcolorbox}

\begin{remark*}[Standard quantified statement]
\[
\operatorname{Field}(F,+,\cdot,0,1)
\Longleftrightarrow
\operatorname{AbelianGroup}(F,+)
\land
\operatorname{AbelianGroup}(F\setminus\{0\},\cdot)
\land
\operatorname{LeftDistributive}(\cdot,+,F)
\land
\operatorname{RightDistributive}(\cdot,+,F).
\]
\end{remark*}

\begin{remark*}[Definition predicate reading]
A field is an additive abelian group whose nonzero elements form a
multiplicative abelian group, with multiplication distributing over addition.
\end{remark*}

\begin{remark*}[Negated quantified statement]
\[
\neg\operatorname{Field}(F,+,\cdot,0,1)
\Longleftrightarrow
\neg\operatorname{AbelianGroup}(F,+)
\lor
\neg\operatorname{AbelianGroup}(F\setminus\{0\},\cdot)
\lor
\neg\operatorname{LeftDistributive}(\cdot,+,F)
\lor
\neg\operatorname{RightDistributive}(\cdot,+,F).
\]
\end{remark*}

\begin{remark*}[Interpretation]
Fields are the algebraic setting where addition, subtraction, multiplication,
and division by nonzero elements are all available.
\end{remark*}

\begin{remark*}[Dependencies]
\begin{itemize}
  \item Abelian group.
  \item Distributivity.
\end{itemize}
\end{remark*}

\begin{proposition}[Every Field is an Integral Domain]
\label{prop:field-is-domain}
Every field is an integral domain.
\hyperref[prf:field-is-domain]{Go to proof.}
\end{proposition}

\begin{remark*}[Standard quantified statement]
\[
\operatorname{Field}(F,+,\cdot,0,1)
\Longrightarrow
\operatorname{IntegralDomain}(F,+,\cdot,0,1).
\]
\end{remark*}

\begin{remark*}[Theorem predicate reading]
The field axioms imply the integral-domain axioms.
\end{remark*}

\begin{remark*}[Negated quantified statement]
\[
\operatorname{Field}(F,+,\cdot,0,1)
\land
\neg\operatorname{IntegralDomain}(F,+,\cdot,0,1)
\Longrightarrow \bot.
\]
\end{remark*}

\begin{remark*}[Interpretation]
The existence of inverses for nonzero elements prevents nonzero zero divisors.
\end{remark*}

\begin{remark*}[Dependencies]
\begin{itemize}
  \item Field.
  \item Integral domain.
  \item Multiplicative inverse.
\end{itemize}
\end{remark*}

\begin{remark*}[Proof TODO]
Proof exists in the archived chapter and should be migrated to a
one-proof-per-file artifact.
Source: \texttt{archive/algebra-and-abstract-structures-legacy.tex}.
\end{remark*}

\begin{remark*}[Examples]
The rational, real, and complex numbers are fields.  The integers are not a
field because nonzero integers such as $2$ do not have multiplicative inverses
inside $\mathbb{Z}$.
\end{remark*}
